
Dự án được tổ chức theo mô hình phân cấp module, mỗi khối chức năng tương ứng với một file Verilog riêng biệt, 
kèm theo các file testbench riêng biệt để kiểm thử từng khối trước khi tích hợp vào hệ thống.

\textbf{Sơ đồ tổ chức thư mục dự án:}
\begin{verbatim}
RISC_CPU/
|
+-- src/                        # Ma nguon thiet ke chinh
|   +-- ALU.v                   # Khoi ALU (8 phep toan, 32-bit)
|   +-- Accumulator.v           # Khoi Accumulator Register
|   +-- program_counter.v       # Khoi Program Counter (PC)
|   +-- address_mux.v           # Khoi Address Mux
|   +-- Instruction_Register.v  # Khoi Instruction Register (IR)
|   +-- memory.v                # Khoi Memory 
|   +-- controller.v            # Khoi Controller
|   +-- CPU_Top.v               # File Top-level tich hop toan he thong
|
+-- tb/                         # Testbench kiem thu tung khoi
|   +-- ALU_tb.v                # Testbench cho ALU
|   +-- Accumulator_tb.v        # Testbench cho Accumulator Register
|   +-- ...                     # Cac testbench tuong ung khac
|   +-- CPU_Top_tb.v            # Testbench toan he thong
|
+-- sim/                        # Ket qua mo phong
    +-- waveform/               # File waveform (.vcd) tu Vivado
\end{verbatim}

\textbf{Vai trò của các thành phần cốt lõi:}

\begin{table}[H]
\centering
\caption{Danh mục vai trò các module trong hệ thống CPU}
\renewcommand{\arraystretch}{1.5}
\begin{tabularx}{\textwidth}{|l|>{\raggedright\arraybackslash}X|}
\hline
\rowcolor[HTML]{EFEFEF} 
\textbf{Tên File} & \textbf{Vai trò trong hệ thống} \\ \hline
\texttt{ALU.v} & Thực thi 8 phép toán số học và logic dựa trên Opcode 3-bit, xuất kết quả 32-bit và tín hiệu \texttt{is\_zero}. \\ \hline
\texttt{Accumulator.v} & Lưu trữ kết quả trung gian từ ALU, hỗ trợ reset đồng bộ và cập nhật dữ liệu khi \texttt{load\_en = 1}. \\ \hline
\texttt{program\_counter.v} & Bộ đếm địa chỉ lệnh 32-bit, hỗ trợ tăng tuần tự hoặc nạp địa chỉ nhảy khi thực hiện lệnh \texttt{JMP}. \\ \hline
\texttt{address\_mux.v} & Lựa chọn giữa địa chỉ lệnh (từ PC) và địa chỉ toán hạng (từ IR) dựa trên tín hiệu điều khiển \texttt{sel}. \\ \hline
\texttt{Instruction\_Register.v} & Lưu trữ mã lệnh nạp từ Memory, thực hiện tách Opcode cho Controller và địa chỉ cho Address Mux. \\ \hline
\texttt{memory.v} & Đơn vị lưu trữ lệnh và dữ liệu, hỗ trợ đọc/ghi thông qua cổng dữ liệu song hướng (\textit{bidirectional port}). \\ \hline
\texttt{controller.v} & Khối điều khiển trung tâm sử dụng FSM 8 trạng thái để phát động các tín hiệu điều khiển toàn hệ thống. \\ \hline
\texttt{CPU\_Top.v} & Module cấp cao nhất thực hiện kết nối tất cả các khối chức năng thành một kiến trúc CPU hoàn chỉnh. \\ \hline
\end{tabularx}
\end{table}